\section{The actual coarse upper rank on each digit class}
\label{sec:rank-upper}

\status{This proves the upper-rank premise retained by SW8 for the
literal block lattices, under \texttt{Std} with $H\ge1$. The integral
path basis and its prefix adaptation are IP3--IP4; their exact dimension
count is CB3. No additional matching or representation premise remains.
The uniform supply of \texttt{Std} is a separate numerical result.
RU9--RU11 also prove the actual segment-to-HH assembly, using RD9's
lower count and RL11's pair lifts.}

\paragraph{RU1--RU3: ranks are weighted path counts.}
Work over a characteristic-zero DVR $O$ with fraction field $K$.
Put $r=X.m$, $k=r+1$ and
\[
 h_r(d)=\binom{d+r-1}{r-1},\qquad
 D_r(\lambda)=\prod_{a<b}
       \frac{\lambda_a-\lambda_b+b-a}{b-a}.                 \tag{RU.1}
\]
For a signed word, a path $\kappa^0=0,\ldots,\kappa^N$ adds a
horizontal strip of the prescribed degree at a vector letter and
removes one at a covector letter. IP4 selects exactly those paths
whose first row obeys every prescribed prefix cutoff. CB3 counts
the terminal standard choices. Thus the actual constrained space is
\[
 \dim_K Y=\sum_{\kappa\ {\rm allowed}}D_r(\kappa^N).
                                                               \tag{RU.2}
\]
The source defines its cutoff as the largest Lie-stable subspace
with the corresponding coordinate-weight bound. In characteristic
zero this is precisely IP4's generic-group bound: exponentiate the
nilpotent root operators and use the diagonal weight decomposition.

For a block at level $(j,j_+)$, SW1's literal multiplication map is
injective: the largest monomial in the independent $S$ variables
recovers that level's coefficient. Its two coefficient factors use
disjoint chain and terminal $\xi$ variables. Clearing denominators
therefore gives
\[
 \operatorname{rank}(\text{block})
   =\sum_{\text{legal }(j,j_+)}
       \dim_KY_{j,j_+}\,h_r(B-j_+).                         \tag{RU.3}
\]
This is also the rank of the actual factorial-weighted block lattice:
it lies in the same rational space and contains a full integral
coefficient lattice.

Two exact counting identities will evaluate the upper bound:
\[
 \begin{aligned}
 \sum_{\kappa\ {\rm unrestricted}}D_r(\kappa^N)
       &=\prod_\ell h_r(d_\ell),\\
 \sum_{\nu_{a+1}\le\lambda_a\le\nu_a}D_r(\lambda)
       &=D_{r+1}(\nu),&
 \sum_{j=0}^{B}h_r(B-j)&=h_{r+1}(B).
 \end{aligned}                                             \tag{RU.4}
\]
The first counts the ordinary monomial basis using IP3. The second
is CB3's tableau branching, valid also for signed dominant weights
after a common determinant shift. The last counts monomials by the
exponent of one distinguished variable. All path sets are finite:
each coordinate is bounded by the total signed degrees.

\paragraph{The actual segment dictionary.}
Fix compatible digits $\rho_i^z,\rho_i^y$ at the $s$ interior
positions and write
\[
 P_i=p^{G-d_i},\quad
 A_i=\frac{b_i-\sum_c\rho^z_{ic}}{P_i},\quad
 C_i=\frac{a_i-\sum_c\rho^y_{ic}}{P_i},\quad
 A_*=\left\lfloor\frac{a-E_0}{p^G}\right\rfloor,\quad B=b_+.
                                                               \tag{RU.5}
\]
Compatibility makes $A_i,C_i$ nonnegative integers.
Split positions into early $i<\sigma$ and late $i\ge\sigma$.
A kept block has fields $(o_i,w_i,s_i,o_+)$, where $s_i$ denotes
the source's \texttt{sg}. Early positions have
$o_i=0$, $w_i\le A_i$, $s_i\le C_i$; late positions have
$o_i\le A_i$, $w_i=s_i=0$. Put
\[
 e_i=\begin{cases}w_i+s_i&i<\sigma,\\o_i&i\ge\sigma,\end{cases}
 \qquad c_n=A_*+\sum_{i<n}(A_i-e_i).                         \tag{RU.6}
\]
The residual word is $V_0,E_0,\ldots,V_{s-1},E_{s-1},U$,
alternating vectors and covectors and ending with a vector. At a
legal level its degrees are
\[
 \begin{array}{c|cc}
 &d_{2i}&d_{2i+1}\\ \hline
 i<\sigma&w_i&C_i-s_i\\
 i\ge\sigma&o_i-j_i&C_i-j_i
 \end{array},
 \qquad d_{2s}=o_+-j_+.                                    \tag{RU.7}
\]
Here $j_i\le\min(o_i,C_i)$ and $j_+\le\min(o_+,B)$.
The actual cutoffs are $\lfloor c_i-M\rfloor$ after each late
$V_i$ and $\lfloor c_s-M\rfloor$ at the endpoint, with $M\ge0$.

\paragraph{RU4--RU5: the missing early row bounds.}
Let $\lambda=\mathrm{La}$, $\epsilon=N_1\lambda/k$,
$\delta=\mathrm{de}$ and $\beta=\mathrm{be}$.
The literal \texttt{Std} inequalities and SW10 give
\[
 \begin{gathered}
 k\ge5,\quad p\ge k^2,\quad s\le N_1,\quad
 40\epsilon\le\delta<1/6,\quad \lambda,N_1\ge1,\\
 0\le\beta<1/12,\qquad
 A_*\ge a/(2Q)\ge k,\qquad A_i(1+\delta)<A_*\quad(i<\sigma).
 \end{gathered}                                             \tag{RU.8}
\]
For $l<i$, increasing physical distances give $P_l\ge pP_i$.
Digit compatibility gives $A_lP_l\le b_l$ and
$b_i-kP_i<A_iP_i$. Using
$b'(1-\beta)\le b_i\le b'(1+\beta)$ yields
$A_l\le2(A_i+k)/p$.
Also \texttt{Kept} gives $s_l\le2\lambda A_*/k$ early.
Consequently
\[
 \sum_{l\le i}(A_l+s_l)
 \le A_i+s\left(\frac{2(A_i+k)}p+\frac{2\lambda A_*}k\right)
 \le A_i+3\epsilon A_*\le A_*.                              \tag{RU.9}
\]
For the last inequalities, $2(A_i+k)/p\le4A_*/k^2
\le\lambda A_*/k$, while
$(A_*-A_i)(1+\delta)>A_*\delta$ gives
$A_*-A_i>34\epsilon A_*$.
Since $w_l\le A_l$, for every $n\le i+1$,
\[
 \sum_{l\le i}w_l-c_n
 \le\sum_{l\le i}(A_l+s_l)-A_*\le0.                         \tag{RU.10}
\]
This proves both neighboring cap bounds without assuming early
caps are monotone.

A vector strip increases the first row by at most its degree;
a covector strip cannot increase it. Hence an early residual path
has $\kappa^{2i+1}_0\le\sum_{l\le i}w_l$.
At late positions the actual cutoff gives
$\kappa^{2i+1}_0\le c_i\le c_{i+1}$.
The terminal cutoff gives $\kappa^{2s+1}_0\le c_s$.
Thus, at every position,
\[
 \kappa^{2i+1}_0\le\min(c_i,c_{i+1}),\qquad
 \kappa^{2s+1}_0\le c_s.                                   \tag{RU.11}
\]

\paragraph{RU6: insert one row and recover all the data.}
Define a path with $r+1$ coordinates by
\[
 \nu^0=0,\quad
 \nu^{2i+1}=(c_i,\kappa^{2i}),\quad
 \nu^{2i+2}=(c_i+A_i-d_{2i},\kappa^{2i+1}).                  \tag{RU.12}
\]
The odd formula includes $i=s$.
It is a path for the unrestricted word
\[
                    +A_*,+A_0,-C_0,\ldots,+A_{s-1},-C_{s-1}.
\]
Indeed, adjoining first entries $f\le f'$ to a horizontal strip
$\alpha/\gamma$ is legal exactly when $\alpha_0\le f$;
its size increases by $f'-f$.
Equation (RU.11) supplies this inequality, and
\[
 c_i+A_i-d_{2i}-c_{i+1}
   =\begin{cases}s_i&i<\sigma,\\j_i&i\ge\sigma,\end{cases}
 \qquad
 c_i+A_i-d_{2i}-c_{i+1}+d_{2i+1}=C_i.                      \tag{RU.13}
\]
The final vector strip and endpoint bound say that
$\kappa^{2s+1}$ interlaces $\nu^{2s+1}$.

Send a source datum to
$(\nu,\kappa^{2s+1},j_+)$.
This preserves its weight
$D_r(\kappa^{2s+1})h_r(B-j_+)$.
It is injective: the last $r$ rows of $\nu$ recover every earlier
$\kappa$, and the stored endpoint recovers the last one.
Strip row-sum differences recover all $d_\ell$. Then early
$(w_i,s_i)=(d_{2i},C_i-d_{2i+1})$; late
$(j_i,o_i)=(C_i-d_{2i+1},d_{2i}+j_i)$; and
$o_+=d_{2s}+j_+$. The fixed digits and zero fields finish recovery.

\paragraph{RU7--RU8: sum the target weights.}
Enlarge the image to all unrestricted target paths, all interlacing
endpoints and $0\le j_+\le B$. Their weights are nonnegative.
Summing in that order with (RU.4), and using (RU.2)--(RU.3), proves
\[
 \boxed{\operatorname{rankOnDigits}
 \le h_{r+1}(A_*)\prod_i h_{r+1}(A_i)h_{r+1}(C_i)
                         h_{r+1}(B)=\texttt{mmB}.}         \tag{RU.14}
\]
This holds for every real $M\ge0$, with no lower-bound margin
restriction. At $M=0$ it is exactly SW8's missing upper inequality.
The numerical companion's uniform implication to \texttt{Std}
therefore supplies the literal \texttt{UpAt} and \texttt{RestUp}
package. Its thresholds may be taken as
$T_0=T_1=1$ and
$T_2=1000(CE_1+CE_2+1)(C_c+1)$, in the source's original order of
constant choices.

\paragraph{RU9: restore the full tensor dimension.}
Here use RD9, under $C_0\ge880$ and \texttt{Std} with $H=4$.
Write $q_j=p^{G_j}$, $q=p^{n'}$, $J=\#\{\text{segments}\}$,
$N=|TT|$, and $R_j=\operatorname{segRank}_j(M)$. Define
\[
 \mathcal B=e^{-\lambda/8}
       +\frac{kQ(M+N_1k)}{b'}+\frac1Y+\frac{kE_0}{a},
 \qquad \kappa_P=C_\kappa N\mathcal B,\quad C_\kappa\ge3.
                                                               \tag{RU.15}
\]
Coordinatewise division $e=\rho+Pu$ gives the exact identity
\[
 \sum_{\substack{0\le\rho_c<P,\ \sum\rho_c\le g\\
                  P\mid g-\sum\rho_c}}
      h_k\!\left(\frac{g-\sum\rho_c}{P}\right)=h_k(g).
                                                               \tag{RU.16}
\]
Summing RD9 over all interior digit rows therefore gives
\[
 R_j\ge L_j h_k(A_{*,j})
       \left(\prod_{i\text{ interior }j}h_k(a_i)h_k(b_i)\right)
       h_k(b_{\rm next}),
 \quad
 L_j=(1-e^{-\lambda/8})(1-e^{-\lambda/3})
              \left(1-\frac{r\lceil M\rceil}{A_{*,j}}\right),
                                                               \tag{RU.17}
\]
when $\lceil M\rceil\le A_{*,j}N_1\lambda/k$.
Kept blocks form a finite set: their digits, $o_i,w_i,s_i$ are
bounded, and $c_s\le A_{*,j}+\sum_i b_i$ bounds the terminal
$o_+$ through its literal kept inequality.

Suppose $\mathcal B<1/6$. Put
$S=kQ/a$, $t=kQ(M+N_1k)/b'$ and
$\delta_f=2/(kY)+(E_0+Q)/a$.
Since $t\ge S(M+5)$ and $A_{*,j}\ge a/(2Q)$,
\[
 \frac{k\lceil M\rceil}{A_{*,j}}<2S(M+1)\le2t<1/3,
 \qquad k\delta_f\le2\mathcal B<1/3.                       \tag{RU.18}
\]
Thus the required margin condition holds. With
$\theta=1-\delta_f$, Bernoulli and the product inequality give
$L_j\theta^r\ge1-3\mathcal B$: the sum of the four losses is at
most
$2e^{-\lambda/8}+S(2M+3)+2/Y+kE_0/a\le3\mathcal B$.
For the actual head degree $g=a_{\rm start}$, its window and the
floor defining $A_{*,j}$ give, for $1\le u\le r$,
\[
 \theta(g+u)\le a-E_0+u-Q\le a-E_0+u-1
                  <q_j(A_{*,j}+u).
\]
Multiplication yields
$\theta^r h_k(g)\le q_j^rh_k(A_{*,j})$.
Hence $(1-3\mathcal B)T_j\le q_j^rR_j$, where
\[
 T_j=h_k(a_{\rm start})\bigl(\prod_{i\text{ interior }j}
                 h_k(a_i)h_k(b_i)\bigr)h_k(b_{\rm next}).
\]

The cyclic intervals partition the selected odd positions, and
the successor permutes their even endpoints. Consequently
$\prod_jT_j=\dim HH$ and $\prod_jq_j=q$.
This is the literal polynomial tensor dimension: CB3 gives
$h_k(d)$ for both vector shape $(d,0,\ldots,0)$ and normalized
covector shape $(d,\ldots,d,0)$; their disjoint-variable tensor
bases multiply. Since $J\le N$,
\[
 \boxed{(1-\kappa_P)\dim HH
       \le\operatorname{gaussian}(q,k)\prod_jR_j.}         \tag{RU.19}
\]
Indeed multiply the segment bounds and use Bernoulli and
$\operatorname{gaussian}(q,k)=1+q+\cdots+q^r\ge q^r$.
If $\kappa_P\ge1$, the claim is immediate; otherwise $N\ge2$
implies $\mathcal B<1/6$. Thus (RU.19) holds for every $M\ge0$.

\paragraph{RU10--RU11: produce the actual HH certificate.}
The weighted block lattice is finite free: in its finite monomial
support the coefficient of $X^\alpha$ lies in
$w^{-v_p(\alpha!)-v_p(\mathrm{specConst})}O$, where $w^2=p$.
Denominator clearing gives its full rational rank. Summing over
the finite kept set proves $\operatorname{rank}_O L_j=R_j$.
Moreover $\mathrm{specConst}\cdot\mathrm{fock}$ is integral on
each block, by adding the two coefficient valuation bounds.

Apply Smith form separately on each block. If a saturated good
submodule $L_{1,j}$ has rank $g_j$, at least $g_j$ diagonal
entries have $v_p\le D_j$. Otherwise its map to those good
coordinates has a nonzero kernel. Divide a kernel vector by its
largest common uniformizer power; saturation retains it and makes
it primitive. Every pairing against $L_j$ then has valuation
$>D_j$, contradicting its good-partner property. The selected
pairs retain single-block support.

For supplied good shares $g_j\ge(1-\kappa_j)R_j$, RL11 (with
$CE_1,CE_2,C_m\ge17$ and $C_t\ge30$) tensors
these choices into actual $F_I\in\mathrm{CC}_e$ and
$G_I\in\mathrm{CC}'_e$ at height $\lfloor t_0\rfloor$, with
\[
 \beta_{HH}(F_I,G_J)=\delta_I\mathbf1_{I=J}+w^c\nu_{IJ},
 \quad 0\ne\delta_I\in O,\quad
 v_p(\delta_I)\le D_*:=\sum_jD_j+4\kappa_0.              \tag{RU.20}
\]
Here $D_j,\kappa_j\ge0$, $c>4D_*$, and the literal normalization
is $\beta_{HH}(F,G)=p^{A_{HH}}G(F)$.
Factor the matrix as $\Delta(I+E)$; all $E_{IJ}$ lie in $wO$
because $v_w(\delta_I)\le2D_*<c$. The adjugate inverts $I+E$
over $O$. If $\delta_{\max}$ has maximal valuation, multiplication
on the right by
\[
 (I+E)^{-1}\operatorname{diag}(\delta_{\max}/\delta_I)
                                                               \tag{RU.21}
\]
gives $\delta_{\max}I$. These are integral column operations,
so they preserve every lattice and precision condition on the
$G_I$, with no height loss. This is the literal
\texttt{FactorsLE} certificate of rank $\prod_jg_j$ and depth $D_*$.

Finally split off $\kappa_P+\sum_j\kappa_j\ge1$, when the desired
mass is nonpositive. Otherwise every remaining factor is nonnegative;
$\prod_j(1-\kappa_j)\ge1-\sum_j\kappa_j$ together with (RU.19)
gives exactly
\[
 \operatorname{RhoGE}_{HH}\!\left(
 D_*,\frac{(1-\kappa_P-\sum_j\kappa_j)\dim HH}
                 {\operatorname{gaussian}(q,k)}\right).  \tag{RU.22}
\]
It is the source's \texttt{aux\_3}/\texttt{HerrExtensionClass}
for the actual $\mathrm{CC}_e,\mathrm{CC}'_e$.
The separate weighted-partner and numerical results supply the
particular $D_j,\kappa_j$ and compare their totals with $D_0,k_x$;
HM then transports this certificate to the harmonic model.
